\frontchapter{Abstract}

Automated ball-launching machines used in sports training operate, to the best of the author's knowledge, in open-loop mode: the trainer configures a fixed angle, speed, and interval, and the machine repeats that programme regardless of the athlete's position or motion. This thesis presents the design, implementation, and quantitative evaluation of a vision-guided ball-launching system that autonomously computes aim parameters for specific human body joints (right knee, right hip, and left shoulder) from real-time multi-view pose estimation in a domestic garage arena. The work is presented as a strong but partial validation: the perception pipeline, the predictive tracker, the safety-gated runtime, and the static and live-aim integration tests are validated end-to-end on real hardware; fully autonomous closed-loop firing at a moving subject remains as the final integration milestone.

The system uses four fixed commodity USB cameras (Hikvision DS-E12, approximately USD~30 each) calibrated using ChArUco boards for intrinsic parameters and AprilTag fiducial markers for extrinsic world-frame registration. Ball detection uses an Ultralytics YOLO detector deployed through TensorRT FP16 with a dynamic-batch engine; pose is estimated using YOLO-Pose with a COCO 17-keypoint output, validated against the MMPose framework. A DLT/SVD triangulator resolves per-camera 2D observations into 3D world-frame coordinates in millimetres, with iterative reprojection-error rejection on the ball pathway and a constant-velocity Kalman filter that smooths the current state and predicts 200--400~ms ahead for the launcher. A ballistic solver computes pitch, yaw, and wheel-motor RPM to direct a Ball Launching Machine (BLM) at the predicted joint position. Low-level actuation is handled by an ESP32 running the \texttt{control\_12\_full.ino} state machine over a 921~600~baud USB-serial link, commanding NEMA-23 stepper motors, a DRV8825-driven pusher, and a counter-rotating flywheel pair. A multi-modal control interface couples the visual targeting loop to an offline speech-recognition front-end via a UDP inter-process channel, supporting hands-free joint selection in training mode.

A six-stage safety-validation protocol governs integration. Stages~S0--S4 were completed on 2026-04-09, including an integrated live test in which the system aimed at and fired upon three named joints (left shoulder, right knee, nose) under operator supervision. The hardware E-STOP latch responds in under 100~ms, and every actuation decision is recorded to a structured JSONL log. A 36-point static ball ground-truth grid yields a corrected mean 3D error of 95.17~mm and 95th-percentile of 166.51~mm; an 81-trial joint-touch protocol (62 valid) yields a mean 3D joint error of 143.38~mm and 95th-percentile of 198.73~mm, both within the per-question acceptance thresholds defined in Chapter~4. The unvalidated boundary is the closed-loop firing regime in which the human target moves during the shot; this is the immediate next experimental milestone.
